\section{The two additive equations in Witt's construction}

The two choices in steps (d) and (e) of the Witt algorithm in
\cite{Camina1997}, pp.~103--104, can be made by explicit finite sums.
The arguments below supply the trace-one argument and the summation
details behind those choices. Actions are written on the left, with
$(st)(x)=s(t(x))$.

\begin{lemma}[Linear independence of distinct field homomorphisms]
\label{lem:field-hom-linear-independence}
\lean{Camina1997.Witt.linearIndependent_smul}
\leanok
Let a group $G$ act faithfully on a field $L$ by field
automorphisms, and let $\varphi_1,\ldots,\varphi_r:L\to L$
be the maps of distinct elements of $G$.
If $a_1,\ldots,a_r\in L$ satisfy
\[
 \sum_{i=1}^r a_i\varphi_i(x)=0\qquad\text{for every }x\in L,
\]
then all the $a_i$ are zero.
\end{lemma}

\begin{proof}
If there is a nonzero relation, choose one with the smallest
positive number of nonzero coefficients. Let $i_0$ have a nonzero
coefficient. There must be a second such index $j$: a relation
supported at $i_0$ alone, evaluated at $1$, would say $a_{i_0}=0$.
Since $\varphi_j\ne\varphi_{i_0}$, choose $b\in F$ with
$\varphi_j(b)\ne\varphi_{i_0}(b)$.
For arbitrary $x$, subtract $\varphi_{i_0}(b)$ times the relation
at $x$ from the relation at $bx$. Multiplicativity gives
\[
 \sum_i a_i\bigl(\varphi_i(b)-\varphi_{i_0}(b)\bigr)
           \varphi_i(x)=0.
\]
The coefficient at $i_0$ is now zero, while the coefficient at
$j$ is nonzero, because it is a product of two nonzero elements
of the field $L$. No previously zero coefficient becomes nonzero.
This contradicts minimality of the chosen support.
\end{proof}

\begin{lemma}[An element of Galois trace one]
\label{lem:galois-trace-one}
\uses{lem:field-hom-linear-independence}
\lean{Camina1997.Witt.exists_sum_smul_eq_one}
\leanok
Let a finite group $G$ act faithfully on a field $L$ by field
automorphisms. There exists $\alpha\in L$ with
\[
 \sum_{s\in G}s(\alpha)=1.
\]
In particular, for a finite Galois extension $L/K$ and its Galois group, the Galois trace $\operatorname{Tr}_{L/K}(x)=\sum_{s\in G}s(x)$
takes the value one, in arbitrary characteristic.
\end{lemma}

\begin{proof}
The distinct automorphisms in the finite group $G$ are linearly
independent as maps $L\to L$ by the preceding lemma. Their sum
is not the zero map: otherwise its coefficients, all equal to
one, would give a nonzero linear relation. Choose $\beta\in L$
such that $z=\sum_{s\in G}s(\beta)\ne0$.
For each $h\in G$, left multiplication by $h$ permutes $G$, so
\[
 h(z)=\sum_{s\in G}(hs)(\beta)=z.
\]
Set $\alpha=\beta/z$. Since every $s$ fixes $z$, we obtain
\[
 \sum_{s\in G}s(\alpha)
   =z^{-1}\sum_{s\in G}s(\beta)=z^{-1}z=1.
\]
In particular this argument does not divide by $|G|$.
\end{proof}

For the following three lemmas, let $L$ be a field of
characteristic $p$, where $p$ is prime, and let a finite group
$G$ act on $L$ by field automorphisms. The prime subfield
$\mathbf F_p\subseteq L$ is fixed by every element of $G$.
Suppose that $\alpha\in L$ satisfies
$\sum_{t\in G}t(\alpha)=1$.
All values in $\mathbf F_p$ in the formulas below are viewed
through this inclusion into $L$.

\begin{lemma}[Solving the additive two-cocycle equation]
\label{lem:witt-two-cocycle-solution}
\source{camina1997:page-0003}
\lean{Camina1997.Witt.twoCocycleSolution_spec, Camina1997.Witt.twoCocycleSolution_one}
\leanok
Let $c:G\times G\to\mathbf F_p$ be a normalized two-cocycle,
meaning that $c(1,s)=c(s,1)=0$ and
\[
 c(s,t)+c(st,u)=c(t,u)+c(s,tu)
 \qquad(s,t,u\in G).
\]
Define
\[
 \delta_s=\sum_{u\in G}c(s,u)\,(su)(\alpha).
\]
Then $\delta_1=0$ and
\[
 \delta_s+s(\delta_t)-\delta_{st}=c(s,t)
 \qquad(s,t\in G).
\]
\end{lemma}

\begin{proof}
The equation $\delta_1=0$ follows term by term from $c(1,u)=0$.
In the sum defining $\delta_s$, reindex by $u\mapsto tu$.
In the sum defining $s(\delta_t)$, use that $s$ fixes
$c(t,u)\in\mathbf F_p$. We get
\[
\begin{aligned}
 \delta_s+s(\delta_t)-\delta_{st}
 &=\sum_{u\in G}
       \bigl(c(s,tu)+c(t,u)-c(st,u)\bigr)(stu)(\alpha)\\
 &=\sum_{u\in G}c(s,t)(stu)(\alpha)\\
 &=c(s,t)\sum_{u\in G}(stu)(\alpha)\\
 &=c(s,t).
\end{aligned}
\]
The second equality is the two-cocycle identity. The final sum
is one because left multiplication by $st$ permutes $G$ and
$\alpha$ has trace one.
\end{proof}

\begin{lemma}[Solving an additive one-cocycle equation]
\label{lem:witt-one-cocycle-solution}
\lean{Camina1997.Witt.oneCocycleSolution_spec}
\leanok
Let $\eta:G\to L$ satisfy
\[
 \eta_{st}=\eta_s+s(\eta_t)\qquad(s,t\in G).
\]
Set
\[
 \gamma=-\sum_{t\in G}\eta_t\,t(\alpha).
\]
Then $s(\gamma)-\gamma=\eta_s$ for every $s\in G$.
\end{lemma}

\begin{proof}
For a fixed $s$, the one-cocycle identity gives
$s(\eta_t)=\eta_{st}-\eta_s$. Therefore
\[
\begin{aligned}
 s(\gamma)
 &=-\sum_{t\in G}s(\eta_t)(st)(\alpha)\\
 &=-\sum_{t\in G}(\eta_{st}-\eta_s)(st)(\alpha)\\
 &=-\sum_{t\in G}\eta_{st}(st)(\alpha)
      +\eta_s\sum_{t\in G}(st)(\alpha)\\
 &=\gamma+\eta_s.
\end{aligned}
\]
Both sums in the last line are obtained from their defining sums
by the permutation $t\mapsto st$, and the second equals one.
Subtracting $\gamma$ gives the assertion.
\end{proof}

\begin{lemma}[The Artin--Schreier image of the two-cocycle solution]
\label{lem:witt-artin-schreier-cocycle}
\uses{lem:witt-two-cocycle-solution,lem:witt-one-cocycle-solution}
\source{camina1997:page-0003,camina1997:page-0004}
\lean{Camina1997.Witt.pow_sub_isOneCocycle_of_twoCocycle, Camina1997.Witt.oneCocycleSolution_spec}
\leanok
For the normalized two-cocycle $c$ and the elements $\delta_s$
above, put $\wp(x)=x^p-x$ and $\eta_s=\wp(\delta_s)$.
Then $\eta_{st}=\eta_s+s(\eta_t)$, and the element
\[
 \gamma=-\sum_{t\in G}\wp(\delta_t)\,t(\alpha)
\]
satisfies
\[
 (s-1)\gamma:=s(\gamma)-\gamma=\wp(\delta_s)
 \qquad(s\in G).
\]
\end{lemma}

\begin{proof}
In characteristic $p$, the binomial theorem gives
$(x+z)^p=x^p+z^p$, since each intermediate binomial coefficient
is divisible by $p$. Thus $\wp$ is additive, preserves negatives,
and commutes with every field automorphism. It vanishes on the
prime subfield, since $c^p=c$ for $c\in\mathbf F_p$.
Apply $\wp$ to
$\delta_s+s(\delta_t)-\delta_{st}=c(s,t)$.
The right side becomes zero, and the left side becomes
$\eta_s+s(\eta_t)-\eta_{st}$, proving the one-cocycle identity.
The preceding lemma applied to this $\eta$ and the given
trace-one element $\alpha$ now proves the formula for $\gamma$.
\end{proof}

\begin{theorem}[Solutions of Witt's two additive equations]
\label{thm:witt-additive-equations}
\uses{lem:galois-trace-one,lem:witt-two-cocycle-solution,lem:witt-artin-schreier-cocycle}
\source{camina1997:page-0003,camina1997:page-0004}
\lean{Camina1997.Witt.exists_witt_equations}
\leanok
Let a finite group $G$ act faithfully on a field $L$ of
characteristic $p$, and let $c:G^2\to\mathbf F_p$ be a normalized
two-cocycle for the trivial action on $\mathbf F_p$.
This includes the Galois action for any finite Galois extension.
There are elements $\delta_s\in L$ and $\gamma\in L$ satisfying
\[
 \delta_s+s(\delta_t)-\delta_{st}=c(s,t),\qquad
 s(\gamma)-\gamma=\delta_s^p-\delta_s
\]
for all $s,t\in G$, with $\delta_1=0$.
\end{theorem}

\begin{proof}
Choose $\alpha$ by the trace-one lemma. Apply the explicit
two-cocycle construction to obtain $\delta_s$, and then apply
the Artin--Schreier cocycle lemma to obtain
$\gamma=-\sum_t\wp(\delta_t)t(\alpha)$.
The two previously proved identities are exactly the required
equations, and the two-cocycle construction also gives
$\delta_1=0$.
\end{proof}
